\SourceLeaf{39}\Article{1}
The principle of virtual motions is best suited to establishing the equation of equilibrium of a system of any number of material points whose motions are subject to any conditions whatsoever. We state it as follows.

Let the system consist of material points $m,m',m''$ etc., at which masses denoted by the same letters are conceived to be concentrated. Let $P$ be one of the accelerating forces acting on point $m$. Imagine that the system is given any infinitely small motion compatible with its conditions, a virtual motion. Let $\mathrm{d}p$ be the motion of point $m$ projected in the direction of the force $P$, that is, multiplied by the cosine of the angle it makes with the direction of force $P$. Finally, let $\Sigma P\mathrm{d}p$ be the sum of all such products for all the forces acting on point $m$. Similarly, let $P'$ stand generally for the forces acting on point $m'$ and $\mathrm{d}p'$ for the motions of point $m'$ projected in the respective directions of those forces, and similarly for the remaining points. With these conventions, the condition for equilibrium of the system is that the sum
\[
m\Sigma P\mathrm{d}p+m'\Sigma P'\mathrm{d}p'+m''\Sigma P''\mathrm{d}p''+\ \text{etc.}
\]
become $=0$ for every virtual motion, as the principle of virtual motions is usually expressed. More precisely, no virtual motion can give this sum a positive value.

\Article{2}
The forces to be considered here fall into three classes.

I. Gravity. Its strength may be taken to be the same at each point and its directions to be parallel. We shall denote its strength by $g$.

\SourceLeaf{40}II. The attractive forces that the points $m,m',m''$ etc. exert on one another. The strength of the attraction is assumed to be proportional to a function of distance, or, equivalently, equal to the product of this function, denoted by the symbol $f$, and the mass concentrated at the attracting point.

III. The forces by which the points $m,m',m''$ etc. are attracted to any number of fixed points. For these forces we shall similarly use the symbol $F$ placed before the distance, and use $M,M',M''$ etc. to denote both the fixed points and the masses assumed to be concentrated at them.

If we now denote the distance between the two points $m,m'$ by $(m,m')$ and likewise denote by $(m,M)$ the distance between the points $m,M$ etc., and if we denote by $z,z',z''$ etc. the heights of the points $m,m',m''$ etc. above an arbitrary horizontal plane $H$, the expression $\Sigma P\mathrm{d}p$ will have the following parts:
\par\begingroup\small
\[
\begin{aligned}
&-g\mathrm{d}z\\
&-m'f(m,m')\mathrm{d}(m,m')-m''f(m,m'')\mathrm{d}(m,m'')-m'''f(m,m''')\mathrm{d}(m,m''')-\ \text{etc.}\\
&-MF(m,M)\mathrm{d}(m,M)-M'F(m,M')\mathrm{d}(m,M')-M''F(m,M'')\mathrm{d}(m,M'')-\ \text{etc.}
\end{aligned}
\]
\endgroup
where the differentials $\mathrm{d}(m,m'),\mathrm{d}(m,m'')$ etc. are partial, since they refer to the virtual motion of point $m$ alone.

In place of the function $f$, let us now introduce the function from which it arises by differentiation. Set $-fx.\mathrm{d}x=\mathrm{d}\varphi x$, or $\int fx.\mathrm{d}x=-\varphi x$. The constant of integration may be chosen at will. If desired, and if the case permits, it may be fixed so that $\varphi\infty=0$, in which case $\varphi t$ will represent the integral $\int fx.\mathrm{d}x$ taken from $x=t$ to $x=\infty$. In just the same way, in place of the function $F$ introduce another function $\Phi$ such that $-Fx.\mathrm{d}x=\mathrm{d}\Phi x$. The expression $\Sigma P\mathrm{d}p$ thus becomes $=$
\[
\begin{aligned}
&-g\mathrm{d}z\\
&+m'\mathrm{d}\varphi(m,m')+m''\mathrm{d}\varphi(m,m'')+m'''\mathrm{d}\varphi(m,m''')+\ \text{etc.}\\
&+M\mathrm{d}\Phi(m,M)+M'\mathrm{d}\Phi(m,M')+M''\mathrm{d}\Phi(m,M'')+\ \text{etc.}
\end{aligned}
\]
Here it should be noted that the differentials in the second line are partial and refer to the motion of point $m$ alone.

But each of these partial differentials clearly has its complement in another expression. Thus both the expression $m\Sigma P\mathrm{d}p$ and the expression $m'\Sigma P'\mathrm{d}p'$ contain the partial differential $mm'\mathrm{d}\varphi(m,m')$. In the first it refers only to the motion of $m$, and in the second only to the motion of $m'$. It follows that the aggre\SourceLeaf{41}gate given in Article 1 is indeed a total differential, specifically $=\mathrm{d}\Omega$, if we set $\Omega=$
\[
\begin{aligned}
&-gmz-gm'z'-gm''z''-\ \text{etc.}\\
&+mm'\varphi(m,m')+mm''\varphi(m,m'')+mm'''\varphi(m,m''')+\ \text{etc.}\\
&\phantom{+mm'\varphi(m,m')}+m'm''\varphi(m',m'')+m'm'''\varphi(m',m''')+\ \text{etc.}\\
&\phantom{+mm'\varphi(m,m')+mm''\varphi(m,m'')}+m''m'''\varphi(m'',m''')+\ \text{etc.}\\
&\phantom{+mm'\varphi(m,m')+mm''\varphi(m,m'')+mm'''\varphi(m,m''')}+\ \text{etc.}\\
&+mM\Phi(m,M)+mM'\Phi(m,M')+mM''\Phi(m,M'')+\ \text{etc.}\\
&+m'M\Phi(m',M)+m'M'\Phi(m',M')+m'M''\Phi(m',M'')+\ \text{etc.}\\
&+m''M\Phi(m'',M)+m''M'\Phi(m'',M')+m''M''\Phi(m'',M'')+\ \text{etc.}\\
&+\ \text{etc.}
\end{aligned}
\]
The condition for equilibrium, then, is that no virtual motion can give the value of the function $\Omega$ a positive increment, or equivalently \textit{that $\Omega$ be a maximum}.

The function $\Omega$ can also be written as follows:
\[
\begin{aligned}
\Omega=\Sigma m\{&-gz+\tfrac12m'\varphi(m,m')+\tfrac12m''\varphi(m,m'')+\tfrac12m'''\varphi(m,m''')+\ \text{etc.}\\
&+M\varphi(m,M)+M'\Phi(m,M')+M''\Phi(m,M'')+\ \text{etc.}\}
\end{aligned}
\]
where the symbol $\Sigma$ represents the sum of the expression written beside it and all the expressions obtained from it when $m$ is exchanged in turn with $m',m'',m'''$ etc.

